% Options for packages loaded elsewhere
\PassOptionsToPackage{unicode}{hyperref}
\PassOptionsToPackage{hyphens}{url}
\PassOptionsToPackage{dvipsnames,svgnames,x11names}{xcolor}
%
\documentclass[
  letterpaper,
  DIV=11,
  numbers=noendperiod]{scrartcl}

\usepackage{amsmath,amssymb}
\usepackage{iftex}
\ifPDFTeX
  \usepackage[T1]{fontenc}
  \usepackage[utf8]{inputenc}
  \usepackage{textcomp} % provide euro and other symbols
\else % if luatex or xetex
  \usepackage{unicode-math}
  \defaultfontfeatures{Scale=MatchLowercase}
  \defaultfontfeatures[\rmfamily]{Ligatures=TeX,Scale=1}
\fi
\usepackage{lmodern}
\ifPDFTeX\else  
    % xetex/luatex font selection
\fi
% Use upquote if available, for straight quotes in verbatim environments
\IfFileExists{upquote.sty}{\usepackage{upquote}}{}
\IfFileExists{microtype.sty}{% use microtype if available
  \usepackage[]{microtype}
  \UseMicrotypeSet[protrusion]{basicmath} % disable protrusion for tt fonts
}{}
\makeatletter
\@ifundefined{KOMAClassName}{% if non-KOMA class
  \IfFileExists{parskip.sty}{%
    \usepackage{parskip}
  }{% else
    \setlength{\parindent}{0pt}
    \setlength{\parskip}{6pt plus 2pt minus 1pt}}
}{% if KOMA class
  \KOMAoptions{parskip=half}}
\makeatother
\usepackage{xcolor}
\setlength{\emergencystretch}{3em} % prevent overfull lines
\setcounter{secnumdepth}{5}
% Make \paragraph and \subparagraph free-standing
\ifx\paragraph\undefined\else
  \let\oldparagraph\paragraph
  \renewcommand{\paragraph}[1]{\oldparagraph{#1}\mbox{}}
\fi
\ifx\subparagraph\undefined\else
  \let\oldsubparagraph\subparagraph
  \renewcommand{\subparagraph}[1]{\oldsubparagraph{#1}\mbox{}}
\fi


\providecommand{\tightlist}{%
  \setlength{\itemsep}{0pt}\setlength{\parskip}{0pt}}\usepackage{longtable,booktabs,array}
\usepackage{calc} % for calculating minipage widths
% Correct order of tables after \paragraph or \subparagraph
\usepackage{etoolbox}
\makeatletter
\patchcmd\longtable{\par}{\if@noskipsec\mbox{}\fi\par}{}{}
\makeatother
% Allow footnotes in longtable head/foot
\IfFileExists{footnotehyper.sty}{\usepackage{footnotehyper}}{\usepackage{footnote}}
\makesavenoteenv{longtable}
\usepackage{graphicx}
\makeatletter
\def\maxwidth{\ifdim\Gin@nat@width>\linewidth\linewidth\else\Gin@nat@width\fi}
\def\maxheight{\ifdim\Gin@nat@height>\textheight\textheight\else\Gin@nat@height\fi}
\makeatother
% Scale images if necessary, so that they will not overflow the page
% margins by default, and it is still possible to overwrite the defaults
% using explicit options in \includegraphics[width, height, ...]{}
\setkeys{Gin}{width=\maxwidth,height=\maxheight,keepaspectratio}
% Set default figure placement to htbp
\makeatletter
\def\fps@figure{htbp}
\makeatother

\usepackage{booktabs}
\usepackage{caption}
\usepackage{longtable}
\usepackage{colortbl}
\usepackage{array}
\usepackage{anyfontsize}
\usepackage{multirow}
\KOMAoption{captions}{tableheading}
\makeatletter
\@ifpackageloaded{caption}{}{\usepackage{caption}}
\AtBeginDocument{%
\ifdefined\contentsname
  \renewcommand*\contentsname{Table of contents}
\else
  \newcommand\contentsname{Table of contents}
\fi
\ifdefined\listfigurename
  \renewcommand*\listfigurename{List of Figures}
\else
  \newcommand\listfigurename{List of Figures}
\fi
\ifdefined\listtablename
  \renewcommand*\listtablename{List of Tables}
\else
  \newcommand\listtablename{List of Tables}
\fi
\ifdefined\figurename
  \renewcommand*\figurename{Figure}
\else
  \newcommand\figurename{Figure}
\fi
\ifdefined\tablename
  \renewcommand*\tablename{Table}
\else
  \newcommand\tablename{Table}
\fi
}
\@ifpackageloaded{float}{}{\usepackage{float}}
\floatstyle{ruled}
\@ifundefined{c@chapter}{\newfloat{codelisting}{h}{lop}}{\newfloat{codelisting}{h}{lop}[chapter]}
\floatname{codelisting}{Listing}
\newcommand*\listoflistings{\listof{codelisting}{List of Listings}}
\makeatother
\makeatletter
\makeatother
\makeatletter
\@ifpackageloaded{caption}{}{\usepackage{caption}}
\@ifpackageloaded{subcaption}{}{\usepackage{subcaption}}
\makeatother
\ifLuaTeX
  \usepackage{selnolig}  % disable illegal ligatures
\fi
\usepackage{bookmark}

\IfFileExists{xurl.sty}{\usepackage{xurl}}{} % add URL line breaks if available
\urlstyle{same} % disable monospaced font for URLs
\hypersetup{
  pdftitle={A Multilevel Bayesian Approach for Inferring the Effect of Fouls on Offensive Patterns in NBA Basketball Games},
  pdfauthor={Michael Light},
  colorlinks=true,
  linkcolor={blue},
  filecolor={Maroon},
  citecolor={Blue},
  urlcolor={Blue},
  pdfcreator={LaTeX via pandoc}}

\title{A Multilevel Bayesian Approach for Inferring the Effect of Fouls
on Offensive Patterns in NBA Basketball Games}
\author{Michael Light}
\date{}

\begin{document}
\maketitle

\section{Timeline}\label{timeline}

Editable
\href{https://docs.google.com/spreadsheets/d/1-yxScU8IQDm0yY-b2s7P6D_-Z3uH5F8GrPLc_FQrKNs/edit?gid=0\#gid=0}{here}.

\begin{table}
\fontsize{12.0pt}{14.4pt}\selectfont
\begin{tabular*}{\linewidth}{@{\extracolsep{\fill}}rl}
\toprule
date & task \\ 
\midrule\addlinespace[2.5pt]
Fri, Apr 18, 2025 & Final Generative Model Validated w/SBC \\ 
Mon, Apr 21, 2025 & Final Computational Model Fit to Real Data \\ 
Wed, Apr 23, 2025 & Initial Draft Ready for Review \\ 
Thu, May 1, 2025 & Final Draft Submitted \\ 
\bottomrule
\end{tabular*}
\end{table}

\section{Project Motivation}\label{project-motivation}

\subsection{Background on the Role of Physicality in NBA
Games}\label{background-on-the-role-of-physicality-in-nba-games}

Basketball is fundamentally a simple game: put a ball through a hoop.
It's also a numbers game: put a ball through a hoop more times than your
opponent. As a result, basketball analysis has historically focused on
offensive aspects - who scored the most points, how they scored them,
who assisted them, and so on. This offensive focus is natural given that
the final score directly reflects which team scored more points.

The introduction of player tracking technology in the early 2010s
briefly revolutionized defensive analysis by providing detailed data on
defender positioning and movement. However, since 2016, this valuable
player tracking data has been made private, creating a significant
challenge for researchers outside NBA organizations.

This data limitation has created a critical gap in basketball analytics.
Without access to detailed player tracking data, researchers are left
with the same fundamental problem they faced before the tracking era:
how can we effectively study and evaluate individual defensive ability
when publicly available, ``counting'' statistics are either:

\begin{itemize}
\item
  Primarily focused on offense
\item
  Only loosely related to defense (e.g., blocks and steals)
\end{itemize}

Solving this problem helps democratize defensive evaluation by creating
metrics that can be derived from publicly available, play-by-play data.
This would allow analysts outside NBA organizations to contribute
meaningful insights to the game and help properly communicate the value
of defensive players who may not show up in traditional box scores.
Further, reliance on player tacking rather than play-by-play data
relegates sample selection to a third-party data provider like
\href{https://sportradar.com/synergy-player-tracking/}{Synergy}, whose
publication aims may introduce bias by selecting observations which may
not be representative of the game as a whole. As such, determining the
viability of using counting stats to infer complex characteristics of
games has the potential to reduce sytematic uncertainty and risk of
future policy changes by third parties.

At a broader level, developing robust defensive metrics from public data
could help bridge the growing divide between basketball analytics and
fan understanding. Currently, there exists a significant gap between
analysts who rely on advanced metrics and fans who primarily use
traditional counting stats like points, rebounds, and assists.
Developing novel, simple, and intuitive ways of understanding the impact
of difficult-to-measure latent player characteristics like defensive
skill represents a crucial challenge in modern basketball analysis.

\subsection{Purpose}\label{purpose}

The purpose of this research is to estimate the causal influence of
individual defender foul counts on the shotmaking probability of
opposing offenses under the hypothesis that players' defensive skill is
a latent trait that can be estimated using NBA play-by-play data.

This work draws widely from fields including sports analytics, causal
inference, Bayesian statistics, and item-response theory.

\subsection{Confounding Expected}\label{confounding-expected}

Among the challenges quantifying any aspect of player skill---offensive,
defensive, or otherwise---is separating individual from team
contribution. For all the simplicity noted in the introduction,
basketball isn't a game you play on paper (Simmons, 2009). We expect any
analysis or model to be imprecise to some degree, and in this research
attempt to account for unobserved confounds associated with shotmaking
probability. These include offensive characteristics including offensive
team, shot location, and shot type.

We narrow our study of defensive skill to offensive posessions that
result in a shot. This means that we define defensive skill primarily as
a component of any defending player's ability to affect the probability
that any opposing offensive player successfully makes an attempted shot.
We do not consider cases where an offensive possession results in a
turnover, such as those caused by steals or shot clock violations, which
although may be indicitative of good or bad defense being played are
ignored for the purpose of this analysis.

As such, any estimate of defensive skill should be interpreted as an
estimate of defensive skill in the context of shotmaking, and should be
assumed independent of higher-level, latent traits like general
defensive skill, general skill, offensive skill, etc.

\subsection{Missing Data}\label{missing-data}

The dataset used for this paper include complete play-by-play records
for every game played during the 2024-2025 NBA season. As our analysis
relies on player-level position data reported on official NBA team
rosters as of 2025-04-14, players who were included in games earlier in
the season but were no longer part of any NBA rosters lacked position
data and were dropped from the sample. This included TK players,
accounting for TK percent of the raw data.

\subsection{Measurement}\label{measurement}

The target of inference (defensive skill) cannot be directly measured,
and the primary measurement (foul count) is measured with some amount of
error resulting from erroneous transcription in raw play\_by\_play data.
Steps taken to account for transcription error can be found in the
supplemental code.

Since foul counts differ between individual defender and defenders tend
to play for fewer minutes as their foul count increases (Chu and Swartz,
2020), the sample size of shots taken in the presence of a defender with
a high foul count versus a low foul count is unbalanced. Multilevel
models and adaptive regularizing priors (McElreath, 2020, pp.~408) are
used to mitigate over fitting such that comparisons can be made across
foul counts.

\section{Literature Review}\label{literature-review}

Bayesian models have been used in a number of previous studies to model
the both offensive and defensive performance. Williams et al.~(2024)
used a Bayesian hierarchical modeling framework to model offensive
performance of teams and players based on their shooting propensities
and abilities. They used discretized shot loactions for the top 100
shot-takers from the 2008/2009--2020/2021 seasons to create an expected
points above average (EPAA) metric. Our project differs in our use of
hierarchical modeling to infer defensive impact, rahter than
contributing to the existing research on offense.

Franks et al.~(2015) was one of the pioneering papers in using Bayesian
hierarchical models to study player tracking data by modeling the
relationship between a defensive assignment and the spacial-temporal
shotmaking tendencies of various offensive players, which is to say how
well they shoot the ball given the location, time, and defender of a
shot attempt. Our project differs in that we use a Bayesian hierarchical
model to infer defensive impact without using raw player tracking data,
rather than modeling the relationship between a known defensive
assignment and shotmaking tendencies.

Finally, Chu and Swartz (2020) use Bayesian methods to model the
time-to-event of a foul as a gamma distributed function of a player,
their accumulated number of fouls, and their position. Their goal was to
develop a decision-making framework for coaches to use in determining
when to sit players at risk of fouling out (usually when a player
reaches 6 fouls) of a game. Although similar in spirit, our project
doesn't model the time-to-event of a foul nor attempt to help coaches
make decisions about sitting their players, but instead focuses on the
effect of fouls while players are on the court.

\section{Modeling Strategy}\label{modeling-strategy}

This section details the steps taken to develop a novel model for
inferring player defensive skill based on the Bayesian workflow proposed
by Gelman and Hill (BDA3, 2006) and others. Naming conventions and
statistical notation are in line with those used by McElreath.

It includes: - a narrative model of the system of inquiry (Betancourt,
TK), - a heuristic causal model motivated by the narrative, - an
agent-based generative model used to simulate sythetic observations of
shots within a basketball game, - a mathematical model, and - a
computational model compatible with the Hamiltonian Monte-Carlo sample
implemented in the Stan probabilistic programming language (Carpenter,
2017).

\subsection{Narrative Model}\label{narrative-model}

We assume that basketball is essentially a game of occupying physical
space. We assume that at any given time T, each location L on a court of
dimension TKxTK has some value V to an offensive player OP, who makes
decisions so as to maximize V. As in chess, defenders are advantaged by
anticipating the spatial preferences (V \textbar{} L) of their opponents
and at time T-1 occupying the location L where V is maximized.

Any given location L can be occupied by no more than one player
(offensive or defensive) at any given time T.

Essentially, at any given time, offensive players prefer to occupy
whatever location on the court maximizes their probability of their team
scoring the basket. Defensive players are able to reduce the value of
any location in one of two ways:

\begin{enumerate}
\def\labelenumi{\arabic{enumi}.}
\tightlist
\item
  By using his defensive skill to preemptively occupy the highest value
  max(V \textbar{} T) location, thereby reducing the total value of any
  other given location
\item
  By using his physical skill to occupy an adjacent location to max(V
  \textbar{} T) and initiating contact at time T, thereby lowering its
  value
\end{enumerate}

This assumes that defensive players' ability to affect shotmaking
probability for any given L is some additive function of his
\textbf{physical skill} and his \textbf{defensive skill}.

We can then assume that defenders with more defensive skill require less
physical skill to affect shotmaking probability for any max(V \textbar{}
T).

We also assume that defenders who use too much physical skill to affect
a given shot are assessed a foul, which accumulate at some rate over the
course of a given game.

We also assume that as a given player's n\_fouls increases, his ability
to express physical skill drops, and so a greater proportion of his
ability to affect shotmaking is comprised of his defensive skill.

For an illustration of the narrative model in action we can think about
the offensive charge, a case in which a defender anticipates the
movements of a driving offensive player to such an accurate degree that
he's able to beat that offender to the location of his shot, resulting
in contact initiated entirely by the offender that is ultimately called
a turnover in favor of the defense. Taking a charge as it's called
(perhaps we should think of it as causing a charge) requires no physical
pressure on the part of the defender, who instead uses his knowledge of
the game and IQ (plausible latent traits of a good defender) to gain the
advantage on a possession. Or, we can think of the cases where a
defender pulls the chair in the post, a rec league classic. Here, a
defender anticipates contact from an offender whose face is away from
the basket, throwing them off balance by moving out of the at the last
moment.

\subsection{Causal model}\label{causal-model}

The complete causal model assumed in this paper can be visualized with
the following Directed Acyclic Graph (DAG):

\includegraphics{report_foul-analysis_files/figure-pdf/unnamed-chunk-2-1.pdf}

This causal diagram includes measurable variables and unmeasurable,
latent player traits which we may infer using an assumed narrative

\section{Generative Model}\label{generative-model}

The complete generative model used for prior elicitation and model
validation via simulation based calibration can be found in the
associated materials.

\begin{verbatim}
# A tibble: 5 x 7
  defensive_team_id defender_id position defender_foul_count
              <int>       <int> <chr>                  <dbl>
1                 1           1 PG                         1
2                 1           1 PG                         1
3                 1           1 PG                         2
4                 1           1 PG                         3
5                 1           1 PG                         3
# i 3 more variables: shot_distance_category <chr>, shots_taken <int>,
#   shots_made <int>
\end{verbatim}

\section{Statistical Models}\label{statistical-models}

The statistical model can be written as the following:

\section{Prior Selection and Predictive
Checks}\label{prior-selection-and-predictive-checks}

\subsection{Two Prior Models for each Model of Foul Effect on
Shot-taking}\label{two-prior-models-for-each-model-of-foul-effect-on-shot-taking}

\subsubsection{Less Constrained}\label{less-constrained}

\subsubsection{More Constrained}\label{more-constrained}

\section{Model Validation}\label{model-validation}

\subsection{Simulation Based
Calibration}\label{simulation-based-calibration}

\section{Observed Sample Selection}\label{observed-sample-selection}

\subsection{Observed Sample of Shots
(Play-by-Play)}\label{observed-sample-of-shots-play-by-play}

\section{Model Fitting and
Diagnostics}\label{model-fitting-and-diagnostics}

\subsection{Fitting Models to Observed
Samples}\label{fitting-models-to-observed-samples}

\subsection{HMC Diagnostics}\label{hmc-diagnostics}

\section{Posterior Analysis}\label{posterior-analysis}

\subsection{Posterior Retrodictive
Checks}\label{posterior-retrodictive-checks}

\subsection{Posterior Residual Checks}\label{posterior-residual-checks}

\section{Counter-factual Analysis}\label{counter-factual-analysis}

\section{Conclusion}\label{conclusion}

\section{Appendices}\label{appendices}

\subsection{Appendix A.}\label{appendix-a.}

\subsubsection{References}\label{references}

Add:

\begin{itemize}
\tightlist
\item
  BDA 3 (Gelman and Hill)
\item
  Statistical rethinking (McElreath)
\item
  Stan (Carpenter, 2017)
\item
  Betancourt narrative model
\end{itemize}

Chu, D., \& Swartz, T. (2020). Foul accumulation in the NBA. Journal of
Quantitative Analysis in Sports, 16(4), 301-309.
https://doi.org/10.1515/jqas-2019-0119

Franks, A., Miller, A., Bornn, L., Goldsberry, K., Characterizing the
spatial structure of defensive skill in professional basketball. (2015)
Ann. Appl. Stat. 9 (1) 94 - 121. https://doi.org/10.1214/14-AOAS799

Simmons, B. (2009). The book of basketball: the NBA according to the
sports guy. Ballantine/ESPN Books.

Williams, B., Schliep, E. M., Fosdick, B., \& Elmore, R. (2024).
Expected points above average: A novel NBA player metric based on
Bayesian hierarchical modeling. arXiv preprint arXiv:2405.10453.
https://arxiv.org/abs/2405.10453



\end{document}
